% GENERATED FILE. Edit canonical lessons, then run npm run generate:books.
% canonical-source-hash: fnv1a64:cb82268f4d77f12a
% canonical-lessons: TE-C510-svartham, TE-C510-nirlaksyam, TE-C510-kramasiksana, TE-C510-atmavisvasam, TE-C510-kutuhalam

\chapter{Habits of Mind}
\label{ch:te-habits-of-mind}
\hlchaptermodality{510}

\begin{chapteropening}
\textbf{By the end of this chapter:} \emph{I can talk about selfishness, carelessness, discipline, confidence and curiosity.}

Complete the last lesson of chapter 510: I can talk about selfishness, carelessness, discipline, confidence and curiosity.
\end{chapteropening}

\section[\texorpdfstring{\te{స్వార్థం} (svārthaṁ)}{స్వార్థం (svārthaṁ)}]{\te{స్వార్థం} (svārthaṁ) — selfishness}

\label{lesson:TE-C510-svartham}



\begin{quote}
Before the new one: say the Telugu for laziness, then the Telugu for high spirits.
\end{quote}

\subsection*{You'll want to know: \te{స్వార్థం}}
\textbf{\te{స్వార్థం}} — \emph{svārthaṁ} — \textquotedblleft{}selfishness\textquotedblright{}.

\begin{quote}
\te{అతను} \te{స్వార్థంతో} \te{అన్నీ} \te{తనకే} \te{తీసుకున్నాడు}.
\end{quote}

\begin{quote}
\emph{atanu svārthantō annī tanakē tīsukunnāḍu.}
\end{quote}

\begin{quote}
Out of selfishness he took everything for himself.
\end{quote}

\te{స్వార్థపరుడు}, a selfish man.

\begin{cousinweb}
\textbf{\te{స్వార్థం}} (\emph{svārthaṁ}) is borrowed from Sanskrit \emph{svārtha}, \textquotedblleft{}one's own affair, self-interest\textquotedblright{} (Monier-Williams), from \emph{sva}, \textquotedblleft{}own\textquotedblright{}, and \emph{artha}, \textquotedblleft{}aim, advantage\textquotedblright{}.
\end{cousinweb}

\begin{grammarlens}[title={What you've built}]
One more word for describing what you lived and felt.
\end{grammarlens}

\subsection*{Your turn}
\begin{itemize}
\raggedright
  \item \emph{Say it:} \emph{svārthaṁ}
  \item \emph{Say it:} \emph{svārthaṁ}, once more
  \item \emph{Say it:} \emph{atanu svārthantō annī tanakē tīsukunnāḍu.} — Out of selfishness he took everything for himself.
  \item \emph{Recall:} say the Telugu for laziness, then the Telugu for high spirits, then say \emph{svārthaṁ} again
\end{itemize}

\begin{tcolorbox}[breakable,colback=teal!4,colframe=teal!35!black,title={Before you move on}]
What does \te{స్వార్థం} mean? (\textquotedblleft{}Selfishness\textquotedblright{}.) Say it once more.
\end{tcolorbox}
\section[\texorpdfstring{\te{నిర్లక్ష్యం} (nirlakṣyaṁ)}{నిర్లక్ష్యం (nirlakṣyaṁ)}]{\te{నిర్లక్ష్యం} (nirlakṣyaṁ) — negligence; neglect}

\label{lesson:TE-C510-nirlaksyam}



\begin{quote}
Before the new one: say the Telugu for high spirits, then the Telugu for selfishness.
\end{quote}

\subsection*{You'll want to know: \te{నిర్లక్ష్యం}}
\textbf{\te{నిర్లక్ష్యం}} — \emph{nirlakṣyaṁ} — \textquotedblleft{}negligence; neglect\textquotedblright{}.

\begin{quote}
\te{డ్రైవరు} \te{నిర్లక్ష్యం} \te{వల్ల} \te{ప్రమాదం} \te{జరిగింది}.
\end{quote}

\begin{quote}
\emph{ḍraivaru nirlakṣyaṁ valla pramādaṁ jarigindi.}
\end{quote}

\begin{quote}
The accident happened because of the driver's negligence.
\end{quote}

\te{నిర్లక్ష్యంగా}, carelessly; -\te{పట్ల} \te{నిర్లక్ష్యం}, neglect of.

\begin{cousinweb}
\textbf{\te{నిర్లక్ష్యం}} (\emph{nirlakṣyaṁ}) is borrowed from Sanskrit \emph{nirlakṣya} (Monier-Williams), from \emph{nir-}, \textquotedblleft{}without\textquotedblright{}, and \emph{lakṣya}, \textquotedblleft{}something to be noticed\textquotedblright{}; Monier-Williams gives \textquotedblleft{}unnoticed\textquotedblright{}, and in Telugu it became the failure to notice: neglect.
\end{cousinweb}

\begin{grammarlens}[title={What you've built}]
One more word for describing what you lived and felt.
\end{grammarlens}

\subsection*{Your turn}
\begin{itemize}
\raggedright
  \item \emph{Say it:} \emph{nirlakṣyaṁ}
  \item \emph{Say it:} \emph{nirlakṣyaṁ}, once more
  \item \emph{Say it:} \emph{ḍraivaru nirlakṣyaṁ valla pramādaṁ jarigindi.} — The accident happened because of the driver's negligence.
  \item \emph{Recall:} say the Telugu for high spirits, then the Telugu for selfishness, then say \emph{nirlakṣyaṁ} again
\end{itemize}

\begin{tcolorbox}[breakable,colback=teal!4,colframe=teal!35!black,title={Before you move on}]
What does \te{నిర్లక్ష్యం} mean? (\textquotedblleft{}Negligence; neglect\textquotedblright{}.) Say it once more.
\end{tcolorbox}
\section[\texorpdfstring{\te{క్రమశిక్షణ} (kramaśikṣaṇa)}{క్రమశిక్షణ (kramaśikṣaṇa)}]{\te{క్రమశిక్షణ} (kramaśikṣaṇa) — discipline}

\label{lesson:TE-C510-kramasiksana}



\begin{quote}
Before the new one: say the Telugu for selfishness, then the Telugu for negligence.
\end{quote}

\subsection*{You'll want to know: \te{క్రమశిక్షణ}}
\textbf{\te{క్రమశిక్షణ}} — \emph{kramaśikṣaṇa} — \textquotedblleft{}discipline\textquotedblright{}.

\begin{quote}
\te{ఆ} \te{బడిలో} \te{క్రమశిక్షణ} \te{చాలా} \te{ముఖ్యం}.
\end{quote}

\begin{quote}
\emph{ā baḍilō kramaśikṣaṇa cālā mukhyaṁ.}
\end{quote}

\begin{quote}
Discipline matters a great deal at that school.
\end{quote}

\te{క్రమం}, order or regularity, + a word for training.

\begin{cousinweb}
\textbf{\te{క్రమశిక్షణ}} (\emph{kramaśikṣaṇa}) is a learned compound built from Sanskrit \emph{krama}, \textquotedblleft{}step, series\textquotedblright{} (CDIAL 3577), and \emph{śikṣaṇa}, \textquotedblleft{}teaching\textquotedblright{} (CDIAL 12429): training in order.
\end{cousinweb}

\begin{grammarlens}[title={What you've built}]
One more word for describing what you lived and felt.
\end{grammarlens}

\subsection*{Your turn}
\begin{itemize}
\raggedright
  \item \emph{Say it:} \emph{kramaśikṣaṇa}
  \item \emph{Say it:} \emph{kramaśikṣaṇa}, once more
  \item \emph{Say it:} \emph{ā baḍilō kramaśikṣaṇa cālā mukhyaṁ.} — Discipline matters a great deal at that school.
  \item \emph{Recall:} say the Telugu for selfishness, then the Telugu for negligence, then say \emph{kramaśikṣaṇa} again
\end{itemize}

\begin{tcolorbox}[breakable,colback=teal!4,colframe=teal!35!black,title={Before you move on}]
What does \te{క్రమశిక్షణ} mean? (\textquotedblleft{}Discipline\textquotedblright{}.) Say it once more.
\end{tcolorbox}
\section[\texorpdfstring{\te{ఆత్మవిశ్వాసం} (ātmaviśvāsaṁ)}{ఆత్మవిశ్వాసం (ātmaviśvāsaṁ)}]{\te{ఆత్మవిశ్వాసం} (ātmaviśvāsaṁ) — self-confidence}

\label{lesson:TE-C510-atmavisvasam}



\begin{quote}
Before the new one: say the Telugu for negligence, then the Telugu for discipline.
\end{quote}

\subsection*{You'll want to know: \te{ఆత్మవిశ్వాసం}}
\textbf{\te{ఆత్మవిశ్వాసం}} — \emph{ātmaviśvāsaṁ} — \textquotedblleft{}self-confidence\textquotedblright{}.

\begin{quote}
\te{మంచి} \te{మార్కులు} \te{రావడంతో} \te{అతనికి} \te{ఆత్మవిశ్వాసం} \te{పెరిగింది}.
\end{quote}

\begin{quote}
\emph{mañci mārkulu rāvaḍantō ataniki ātmaviśvāsaṁ perigindi.}
\end{quote}

\begin{quote}
With good marks his self-confidence grew.
\end{quote}

\te{ఆత్మ}, self, + \te{విశ్వాసం}, trust: trust in yourself.

\begin{cousinweb}
\textbf{\te{ఆత్మవిశ్వాసం}} (\emph{ātmaviśvāsaṁ}) is a learned compound built from Sanskrit \emph{ātman}, \textquotedblleft{}breath, self\textquotedblright{} (CDIAL 1135), and \emph{viśvāsa}, \textquotedblleft{}trust, confidence\textquotedblright{} (CDIAL 11966).
\end{cousinweb}

\begin{grammarlens}[title={What you've built}]
One more word for describing what you lived and felt.
\end{grammarlens}

\subsection*{Your turn}
\begin{itemize}
\raggedright
  \item \emph{Say it:} \emph{ātmaviśvāsaṁ}
  \item \emph{Say it:} \emph{ātmaviśvāsaṁ}, once more
  \item \emph{Say it:} \emph{mañci mārkulu rāvaḍantō ataniki ātmaviśvāsaṁ perigindi.} — With good marks his self-confidence grew.
  \item \emph{Recall:} say the Telugu for negligence, then the Telugu for discipline, then say \emph{ātmaviśvāsaṁ} again
\end{itemize}

\begin{tcolorbox}[breakable,colback=teal!4,colframe=teal!35!black,title={Before you move on}]
What does \te{ఆత్మవిశ్వాసం} mean? (\textquotedblleft{}Self-confidence\textquotedblright{}.) Say it once more.
\end{tcolorbox}
\section[\texorpdfstring{\te{కుతూహలం} (kutūhalaṁ)}{కుతూహలం (kutūhalaṁ)}]{\te{కుతూహలం} (kutūhalaṁ) — curiosity}

\label{lesson:TE-C510-kutuhalam}



\begin{quote}
Before the new one: say the Telugu for discipline, then the Telugu for self-confidence.
\end{quote}

\subsection*{You'll want to know: \te{కుతూహలం}}
\textbf{\te{కుతూహలం}} — \emph{kutūhalaṁ} — \textquotedblleft{}curiosity\textquotedblright{}.

\begin{quote}
\te{పిల్లలు} \te{కుతూహలంతో} \te{పెట్టె} \te{తెరిచారు}.
\end{quote}

\begin{quote}
\emph{pillalu kutūhalantō peṭṭe tericāru.}
\end{quote}

\begin{quote}
The children opened the box with curiosity.
\end{quote}

\te{కుతూహలంగా}, curiously, eagerly.

\begin{cousinweb}
\textbf{\te{కుతూహలం}} (\emph{kutūhalaṁ}) is borrowed from Sanskrit \emph{kutūhala}, \textquotedblleft{}curiosity, interest in anything unusual\textquotedblright{} (Monier-Williams).
\end{cousinweb}

\begin{grammarlens}[title={What you've built}]
One more word for describing what you lived and felt.
\end{grammarlens}

\subsection*{Your turn}
\begin{itemize}
\raggedright
  \item \emph{Say it:} \emph{kutūhalaṁ}
  \item \emph{Say it:} \emph{kutūhalaṁ}, once more
  \item \emph{Say it:} \emph{pillalu kutūhalantō peṭṭe tericāru.} — The children opened the box with curiosity.
  \item \emph{Recall:} say the Telugu for discipline, then the Telugu for self-confidence, then say \emph{kutūhalaṁ} again
\end{itemize}

\begin{tcolorbox}[breakable,colback=teal!4,colframe=teal!35!black,title={Before you move on}]
What does \te{కుతూహలం} mean? (\textquotedblleft{}Curiosity\textquotedblright{}.) Say it once more.
\end{tcolorbox}
